% !TEX program = pdflatex
%=====================================================================
%  Criptografía y Seguridad - ITBA
%  Clase 11 — Protección de datos
%  Apunte integral basado en la teoría (transcript + diapositivas)
%=====================================================================
\documentclass[11pt,a4paper]{article}

%---------------------------------------------------------------------
%  Paquetes
%---------------------------------------------------------------------
\usepackage[utf8]{inputenc}
\usepackage[T1]{fontenc}
\usepackage[spanish,es-tabla,es-noshorthands]{babel}
\usepackage{lmodern}
\usepackage{microtype}

\usepackage{amsmath,amssymb,amsthm}
\usepackage{mathtools}
\usepackage{geometry}
\geometry{a4paper,margin=2.4cm,headheight=22pt}

\usepackage{xcolor}
\usepackage{graphicx}
\usepackage{enumitem}
\usepackage{booktabs}
\usepackage{array}
\usepackage{tcolorbox}
\tcbuselibrary{skins,breakable}
\usepackage{fancyhdr}
\usepackage{titlesec}
\usepackage{hyperref}
\usepackage{listings}
\usepackage{caption}

\usepackage{tikz}
\usetikzlibrary{shapes.geometric,shapes.misc,shapes.symbols,arrows.meta,positioning,calc,
                fit,backgrounds,decorations.pathmorphing,shadows.blur,chains}

%---------------------------------------------------------------------
%  Paleta de color (inspirada en las diapositivas de la cátedra)
%---------------------------------------------------------------------
\definecolor{itbaCyan}{HTML}{6EC6D9}
\definecolor{itbaCyanDark}{HTML}{2E8B9E}
\definecolor{itbaBlue}{HTML}{AEDCEA}
\definecolor{boxBlue}{HTML}{BBD4F0}
\definecolor{slideGray}{HTML}{ECECEC}
\definecolor{practGreen}{HTML}{4F8A3D}
\definecolor{practGreenL}{HTML}{E2EFD9}
\definecolor{attackRed}{HTML}{C0392B}
\definecolor{codeBg}{HTML}{F5F5F5}
% Colores auxiliares para los estados del dato (idénticos al diagrama de la slide)
\definecolor{stateRest}{HTML}{F4B6B0}   % rosa  - at rest
\definecolor{stateMotion}{HTML}{F4CE5E} % amarillo - in motion
\definecolor{stateUse}{HTML}{CBC6E6}    % lila  - in use

%---------------------------------------------------------------------
%  Hipervínculos
%---------------------------------------------------------------------
\hypersetup{
  colorlinks=true,
  linkcolor=itbaCyanDark,
  urlcolor=itbaCyanDark,
  citecolor=itbaCyanDark,
  pdftitle={Clase 11 - Proteccion de datos},
  pdfauthor={Criptografia y Seguridad - ITBA}
}

%---------------------------------------------------------------------
%  Encabezados y pies
%---------------------------------------------------------------------
\pagestyle{fancy}
\fancyhf{}
\lhead{\small\textsc{Criptografía y Seguridad — ITBA}}
\rhead{\small Clase 11 · Protección de datos}
\cfoot{\small\thepage}
\renewcommand{\headrulewidth}{0.4pt}
\renewcommand{\headrule}{\hbox to\headwidth{\color{itbaCyanDark}\leaders\hrule height \headrulewidth\hfill}}

%---------------------------------------------------------------------
%  Títulos de sección con barra de color (estilo diapositiva)
%---------------------------------------------------------------------
\titleformat{\section}
  {\normalfont\Large\bfseries\color{itbaCyanDark}}
  {\thesection}{0.6em}{}
  [\vspace{-0.4em}{\color{itbaCyan}\titlerule[2pt]}]
\titleformat{\subsection}
  {\normalfont\large\bfseries\color{itbaCyanDark}}
  {\thesubsection}{0.6em}{}
\titleformat{\subsubsection}
  {\normalfont\normalsize\bfseries\color{black}}
  {\thesubsubsection}{0.6em}{}

%---------------------------------------------------------------------
%  Cajas reutilizables
%---------------------------------------------------------------------
% Caja "diapositiva" (recuadro gris de las slides)
\newtcolorbox{slidebox}[1][]{
  enhanced, colback=slideGray, colframe=black!55, boxrule=0.5pt,
  arc=1pt, left=6pt,right=6pt,top=4pt,bottom=4pt, #1}

% Caja de definición
\newtcolorbox{defbox}[1]{
  enhanced, breakable, colback=itbaBlue!25, colframe=itbaCyanDark,
  boxrule=1pt, arc=2pt, left=8pt,right=8pt,top=6pt,bottom=6pt,
  fonttitle=\bfseries, coltitle=white, title={#1}}

% Caja de nota / aclaración del profe (verde)
\newtcolorbox{notebox}[1]{
  enhanced, breakable, colback=practGreenL, colframe=practGreen,
  boxrule=1pt, arc=2pt, left=8pt,right=8pt,top=6pt,bottom=6pt,
  fonttitle=\bfseries, coltitle=white, title={#1}}

% Caja de atención / ataque / error
\newtcolorbox{warnbox}[1]{
  enhanced, breakable, colback=attackRed!8, colframe=attackRed,
  boxrule=1pt, arc=2pt, left=8pt,right=8pt,top=6pt,bottom=6pt,
  fonttitle=\bfseries, coltitle=white, title={#1}}

% Caja "recuadro rojo" de las diapositivas (frase destacada en rojo)
\newtcolorbox{redslide}{
  enhanced, colback=attackRed!85, colframe=attackRed!85, boxrule=0pt,
  arc=2pt, halign=center, coltext=white, fontupper=\large\bfseries,
  left=8pt,right=8pt,top=8pt,bottom=8pt}

%---------------------------------------------------------------------
%  Estilo de código
%---------------------------------------------------------------------
\lstdefinestyle{json}{
  basicstyle=\ttfamily\footnotesize,
  backgroundcolor=\color{codeBg},
  keywordstyle=\color{itbaCyanDark}\bfseries,
  commentstyle=\color{practGreen}\itshape,
  stringstyle=\color{black},
  breaklines=true, frame=single, rulecolor=\color{black!30},
  showstringspaces=false, columns=fullflexible, tabsize=2,
  xleftmargin=10pt, framexleftmargin=8pt}

%---------------------------------------------------------------------
%  Cita textual del profe
%---------------------------------------------------------------------
\newcommand{\prof}[1]{%
  \medskip\noindent\hspace*{1em}\textit{``#1''}\\[-1pt]%
  \hspace*{1em}{\footnotesize\textcolor{black!55}{--- Prof.\ Leandro Suárez (clase teórica)}}\medskip}

%=====================================================================
\begin{document}
%=====================================================================

%---------------------------------------------------------------------
%  Portada
%---------------------------------------------------------------------
\begin{titlepage}
\centering
\vspace*{1.0cm}

\begin{tikzpicture}
  \node[rounded corners=2pt, fill=itbaCyan, text=black, inner sep=14pt,
        minimum width=0.8\textwidth, align=center] (t)
        {\Huge\bfseries Criptografía y Seguridad};
\end{tikzpicture}

\vspace{0.6cm}
{\LARGE Seguridad en la empresa}\\[0.25cm]
{\Huge\bfseries\color{itbaCyanDark} Protección de datos}

\vspace{0.8cm}

% Ilustración: bóveda / caja fuerte (recreación del icono de portada)
\begin{tikzpicture}[scale=1.0]
  \fill[black!12] (-3.2,-2.4) rectangle (3.2,2.8);
  \draw[line width=2pt, black!55] (-3.2,-2.4) rectangle (3.2,2.8);
  \fill[itbaCyanDark!75] (-2.4,-2.0) rectangle (2.4,2.4);
  \draw[line width=3pt, black!70] (-2.4,-2.0) rectangle (2.4,2.4);
  \draw[line width=1.2pt, itbaCyan] (-2.0,-1.6) rectangle (2.0,2.0);
  \draw[line width=2pt, black!70] (0,-2.0) -- (0,2.4);
  \foreach \a in {0,45,...,315}{
    \draw[line width=2pt, black!70] (-1.1,0.2) -- ++(\a:0.6);
  }
  \draw[line width=2.5pt, black!70] (-1.1,0.2) circle (0.42);
  \fill[itbaCyan] (-1.1,0.2) circle (0.16);
  \draw[line width=2.5pt, black!70] (1.1,0.2) circle (0.42);
  \fill[itbaCyan] (1.1,0.2) circle (0.16);
  \foreach \y in {-1.2,-0.4,0.8,1.6}{
     \fill[black!70] (-2.2,\y) circle (0.07);
     \fill[black!70] (2.2,\y) circle (0.07);
  }
\end{tikzpicture}

\vfill
{\large Apunte integral — Teoría}\\[0.2cm]
{\normalsize Clase 11 \quad·\quad Última clase de la materia}\\[0.4cm]
{\small Instituto Tecnológico de Buenos Aires (ITBA)}

\vspace{0.6cm}
{\footnotesize\itshape Documento elaborado a partir de la transcripción de la clase
teórica (Prof.\ Leandro Suárez) y de las diapositivas oficiales. No se agregó contenido
externo al material de clase.}
\end{titlepage}

\tableofcontents
\newpage

%=====================================================================
\section{Introducción: el problema}
%=====================================================================

Esta es la \textbf{última clase de la materia} y conecta directamente con la anterior
(\emph{Seguridad en la empresa}). El tema se mira desde la perspectiva de un
\textbf{CISO} (\emph{Chief Information Security Officer}): el director de seguridad de la
organización, la persona encargada de velar por la seguridad de la empresa y que,
generalmente, tiene \textbf{implicancias legales} si se produjera un incidente de
seguridad severo con consecuencias jurídicas.

\prof{Acá es donde pasamos a aplicar, digamos, para la protección de datos se pasa a aplicar
todos los principios que ya vimos en pos de proteger la CIA: la confidencialidad, integridad
y disponibilidad de los datos.}

\subsection{El escenario que motiva todo}

\begin{slidebox}
En una empresa se define una \textbf{política de protección de datos personales} que dice:
\emph{«Un empleado no puede acceder a ver los sueldos de otros empleados.»}
\end{slidebox}

El profesor plantea el problema como disparador: \emph{¿dónde habría que poner controles
para proteger esa información?} La respuesta no es trivial, porque el mismo dato (la planilla
de sueldos) se replica y viaja por muchísimos lugares. La diapositiva muestra un diagrama
de flujo del dato que conviene tener presente durante toda la clase: un dato «sensible» no
vive en un solo lugar.

\begin{center}
\begin{tikzpicture}[
  font=\footnotesize,
  store/.style={cylinder,shape border rotate=90,draw,minimum width=1.5cm,minimum height=0.9cm,align=center},
  app/.style={rounded corners=6pt,draw,minimum width=1.6cm,minimum height=1.0cm,align=center},
  doc/.style={draw,signal,signal to=south,minimum width=1.5cm,minimum height=0.9cm,align=center,inner sep=2pt},
  cl/.style={cloud,draw,aspect=2,minimum width=1.7cm,minimum height=1.1cm,align=center},
  >=Stealth, node distance=0.9cm and 0.95cm]

  \node[store] (rh) {RR.HH.};
  \node[app, right=of rh] (app) {RR.HH.\\APP};
  \node[cl, below right=0.9cm and 0.2cm of app] (drive) {Drive};
  \node[app, right=1.7cm of app] (xls) {Excel};
  \node[app, right=of xls] (mail) {Email};
  \node[cl, right=of mail] (nube) {};

  \node[doc, below=0.5cm of rh] (b1) {BACKUP};
  \node[doc, below=0.45cm of drive] (b2) {BACKUP};

  \draw[->] (rh) -- (app);
  \draw[->] (rh) -- (b1);
  \draw[->] (app) |- (drive);
  \draw[->] (drive) -- (b2);
  \draw[->] (app.east) -- ([yshift=4pt]xls.west);
  \draw[->] (drive.east) -| (xls.south);
  \draw[->] (xls) -- (mail);
  \draw[->] (mail) -- (nube);
\end{tikzpicture}
\end{center}
{\centering\footnotesize El mismo dato sensible se copia en una base, una app, backups,
una hoja de cálculo, un Drive, un mail y la nube.\par}

\begin{warnbox}{Idea central de la clase}
Los \textbf{controles de acceso} de un sistema son parte de la solución, \textbf{pero no
son suficientes}. El dato escapa del sistema (se exporta a un Excel, se sube a Drive, se
manda por mail, se respalda en backups), y en cada uno de esos lugares hay que protegerlo.
Por eso se necesita un enfoque integral: el \emph{ciclo de protección de la información}.
\end{warnbox}

%=====================================================================
\section{El ciclo de protección de la información}
%=====================================================================

La materia vuelve sobre la idea de los \textbf{ciclos} en seguridad. La protección de la
información se organiza como un ciclo de seis etapas que se repite indefinidamente.

\begin{center}
\begin{tikzpicture}[
  font=\small,
  stage/.style={draw, rounded corners=2pt, fill=white, minimum width=2.6cm,
                minimum height=1.0cm, align=center, line width=0.8pt},
  >=Stealth]

  \def\R{3.5cm}
  \node[stage] (gov)  at (90:\R)  {Gobierno\\\footnotesize(Govern)};
  \node[stage] (dis)  at (30:\R)  {Descubrimiento\\\footnotesize(Discover)};
  \node[stage] (pro)  at (-30:\R) {Protección\\\footnotesize(Protect)};
  \node[stage] (com)  at (-90:\R) {Cumplimiento\\\footnotesize(Comply)};
  \node[stage] (det)  at (-150:\R){Detección\\\footnotesize(Detect)};
  \node[stage] (res)  at (150:\R) {Respuesta\\\footnotesize(Respond)};

  \draw[->,line width=1pt] (gov) to[bend left=15] (dis);
  \draw[->,line width=1pt] (dis) to[bend left=15] (pro);
  \draw[->,line width=1pt] (pro) to[bend left=15] (com);
  \draw[->,line width=1pt] (com) to[bend left=15] (det);
  \draw[->,line width=1pt] (det) to[bend left=15] (res);
  \draw[->,line width=1pt] (res) to[bend left=15] (gov);
\end{tikzpicture}
\end{center}

\begin{enumerate}[leftmargin=1.6em,itemsep=3pt]
  \item \textbf{Gobierno (Govern):} el punto de \emph{gestión}, donde se toman las
        decisiones y se define cómo se ejecuta todo. (Es lo visto en la clase anterior,
        \emph{Seguridad en la empresa}.)
  \item \textbf{Descubrimiento (Discover):} ver qué datos hay en el universo de datos de
        la organización.
  \item \textbf{Protección (Protect):} definir cómo y qué datos se protegen (no todos los
        datos son sensibles).
  \item \textbf{Cumplimiento (Comply):} asegurar que esa protección se cumple efectivamente.
  \item \textbf{Detección (Detect):} ser capaces de detectar filtraciones o brechas de la
        información.
  \item \textbf{Respuesta (Respond):} responder a esas filtraciones y evitar que se sigan
        produciendo.
\end{enumerate}

\noindent Y luego se vuelve a empezar: \prof{Y luego volvemos a empezar. Es un ciclo.}

\noindent El resto del apunte recorre el ciclo etapa por etapa.

%=====================================================================
\section{Gobierno (Data Governance)}
%=====================================================================

\begin{defbox}{Data Governance}
Son \textbf{todas las tareas que se realizan para asegurar que los datos sean seguros,
estén íntegros, sean confidenciales y estén disponibles} en todo momento que se requiera.
Incluye:
\begin{itemize}[leftmargin=1.4em,itemsep=1pt,topsep=2pt]
  \item las medidas que deben tomar las \textbf{personas},
  \item los \textbf{procesos} que deben seguir,
  \item la \textbf{tecnología} que los respalda,
\end{itemize}
durante todo el \textbf{ciclo de vida de los datos}. También \textbf{define los roles}
dentro de la organización para el uso de la información.
\end{defbox}

Como los sistemas, los datos también tienen un ciclo de vida. Hay normativas que obligan a
que los datos vivan cierta cantidad de tiempo, que se borren o que no se borren; y casos en
los que se deben eliminar datos por cuestiones legales.

El gobierno de datos establece \textbf{estándares internos} (conocidos como
\textbf{políticas de datos}) que se aplican a cinco actividades sobre los datos:

\begin{center}
\begin{tikzpicture}[font=\small,>=Stealth,
  dgstep/.style={draw,rounded corners=2pt,fill=itbaBlue!30,minimum height=0.9cm,
               minimum width=2.5cm,align=center}]
  \node[dgstep] (a) {Recopilación};
  \node[dgstep,right=0.4cm of a] (b) {Almacenamiento};
  \node[dgstep,right=0.4cm of b] (c) {Procesamiento};
  \node[dgstep,below=0.5cm of b] (d) {Retención};
  \node[dgstep,right=0.4cm of d] (e) {Eliminación};
\end{tikzpicture}
\end{center}

\subsection{El ciclo de vida del dato (Data life cycle)}

La diapositiva lo presenta como un ciclo de 5 fases:

\begin{enumerate}[leftmargin=1.8em,itemsep=2pt]
  \item \textbf{Create \& collect} — recolección de información a través de todos los
        canales disponibles.
  \item \textbf{Store \& transmit} — almacenamiento y transmisión de lo recolectado para su
        procesamiento (a través de sistemas IT, archivado físico y transporte).
  \item \textbf{Use \& distribute} — procesamiento de la información y distribución interna.
  \item \textbf{Retain} — retención de la información procesada por posibles usos y
        requerimientos legales/regulatorios.
  \item \textbf{Dispose and destroy} — disposición/destrucción de la información una vez que
        vence el período de retención.
\end{enumerate}

\begin{notebox}{¿Por qué importa el ciclo de vida?}
\prof{Tener datos en el fondo es tener información, es tener activos importantes que buscan
los atacantes muchas veces.} Por eso conviene gestionar cuándo y cuánto viven los datos:
cuanto más datos sensibles se conservan sin necesidad, mayor es la superficie de ataque.
\end{notebox}

%=====================================================================
\section{Regulaciones (externas e internacionales)}
%=====================================================================

El gobierno de datos está condicionado por regulaciones que la organización debe cumplir.

\subsection{Regulación local: Ley 25.326 (Argentina)}

La \textbf{Ley 25.326 de Protección de Datos Personales} aplica a todo sistema que opera en
Argentina. La clase resaltó dos artículos:

\begin{slidebox}
\textbf{Artículo 1º (Objeto).} «La presente ley tiene por objeto la protección integral de
los datos personales asentados en archivos, registros, bancos de datos, u otros medios
técnicos de tratamiento de datos, sean éstos públicos, o privados destinados a dar informes,
para garantizar el derecho al honor y a la intimidad de las personas, así como también el
acceso a la información que sobre las mismas se registre, de conformidad a lo establecido en
el artículo 43, párrafo tercero de la Constitución Nacional.»

\smallskip
\textbf{Artículo 9º (Seguridad de los datos).}
\emph{1.} «El responsable o usuario del archivo de datos debe adoptar las medidas técnicas y
organizativas que resulten necesarias para garantizar la seguridad y confidencialidad de los
datos personales, de modo de evitar su adulteración, pérdida, consulta o tratamiento no
autorizado, y que permitan detectar desviaciones, intencionales o no, de información, ya sea
que los riesgos provengan de la acción humana o del medio técnico utilizado.»
\emph{2.} «Queda prohibido registrar datos personales en archivos, registros o bancos que no
reúnan condiciones técnicas de integridad y seguridad.»
\end{slidebox}

\begin{defbox}{¿Qué es un dato personal?}
Es todo lo que permite \textbf{identificar a una persona} sólo por ese mismo dato. Ejemplos:
nombre y apellido, dirección de correo, número de documento, dirección del domicilio.

\smallskip
El nombre es un caso de borde: un nombre muy común (muchos «Leandro Suárez») casi no
identifica, mientras que uno único (p.\ ej.\ «Lionel Messi») sí. Igualmente, las leyes lo
consideran dato personal por el caso particular en que sí es identificable.
\end{defbox}

\begin{notebox}{Lo importante de esta ley (según el profe)}
La ley \textbf{no especifica una tecnología}: te dice «a grandes rasgos, hacé todo lo posible
para que estos datos permanezcan íntegros y confidenciales y evitá filtraciones o borrados no
intencionales». Esto es bueno: al no atarse a una tecnología (la ley es de los años 80),
\textbf{sigue siendo aplicable hoy}.

\smallskip
Consecuencia práctica: si tu aplicación guarda datos personales, hay que \textbf{asesorarse
legalmente}. Y este tipo de leyes suele dar a las personas el derecho a \emph{ver} qué datos
guarda la organización, \emph{modificarlos} e incluso \emph{pedir su borrado} (el
«\textbf{derecho al olvido}»). Es algo que se exige en casi todo el mundo (Facebook, Twitter,
etc., permiten descargar tu historial y borrar tu cuenta).
\end{notebox}

\subsection{Regulaciones internacionales}

\begin{center}
\renewcommand{\arraystretch}{1.35}
\begin{tabular}{>{\bfseries}p{2.4cm} p{12.2cm}}
\toprule
HIPAA & \emph{Health Insurance Portability and Accountability Act of 1996.} Protección de
datos relacionados con la salud de las personas (\textbf{PHI / PPI}). Surge de acá la
\textbf{anonimización} de los datos médicos (ver más abajo). \\
GDPR & \emph{General Data Protection Regulation.} La más conocida de protección de datos
personales; de la Unión Europea, aplicada en \textbf{2018}. \\
PCI-DSS & \emph{Payment Card Industry Data Security Standard.} Regula el uso de información de
tarjetas de crédito. \textbf{Obligatorio} para cualquier empresa que quiera operar con
tarjetas. \\
SOX & \emph{Sarbanes-Oxley Act} (ley federal de EE.UU., 2002). Para empresas que cotizan en
bolsa / públicas; exige transparencia en sus operaciones. \\
FFIEC & \emph{Assessment Tool.} Framework (hoy algo en desuso). \\
FDA & \emph{Cybersecurity.} Consideraciones de ciberseguridad de la \emph{Food and Drug
Administration} (estudios médicos, alimentos y drogas). \\
\bottomrule
\end{tabular}
\end{center}

\paragraph{GDPR y las cookies.} Por esta ley (2018) apareció «de golpe» el tema de los
banners de cookies: GDPR pide un \textbf{consentimiento explícito} del usuario. Por eso los
sitios te dejan elegir «solo lo esencial» (lo necesario para que el sitio funcione) y
rechazar las demás (marketing y similares, que registran tus preferencias/actividad e incluso
pueden acceder a cookies de terceros). Aplica aunque el sistema esté en Argentina, siempre
que haya usuarios o infraestructura en la Unión Europea.

\begin{notebox}{PCI-DSS: dos conceptos clave que impuso}
\textbf{1) Acotar la cantidad de datos a proteger según su criticidad.} PCI definió que los
datos críticos de una tarjeta son básicamente \textbf{5}: número de tarjeta, nombre del
titular, vencimiento, código de seguridad (el más crítico) y datos de la banda magnética/chip
(a veces, un \emph{PIN} que funciona como segundo factor en el POS). Al acotar a esos
\(\sim\)5 datos, la norma da un \emph{blueprint}: concentrarse en proteger lo realmente
sensible y no «todo».
\prof{Una de las cosas más importantes a la hora de proteger datos es acotar la cantidad de
datos en base a su criticidad.}

\medskip
\textbf{2) Tokenización.} La tarjeta (sus 5 datos) se guarda «bajo 7 llaves» en una base muy
protegida y segmentada; el sistema opera con un \textbf{token} que apunta a esos datos. Si se
filtrara el token, no tiene valor por sí mismo: habría que ir a buscar los datos a la base
recontra-protegida. Así se desacopla la operación de los datos sensibles.
\end{notebox}

\begin{notebox}{HIPAA: anonimización}
Con HIPAA surgió el concepto de \textbf{anonimizar} los datos médicos. La idea: se pueden
seguir usando los datos para análisis y búsqueda de patrones, pero \emph{sin saber de quién
es cada dato}. \prof{Yo no sé de quién es este estudio, yo no sé de quién es esta enfermedad.}
Esto se suele exigir para historias clínicas.
\end{notebox}

%=====================================================================
\section{Data Roles (roles sobre los datos)}
%=====================================================================

Para distribuir la responsabilidad sobre los datos se define una \textbf{jerarquía de roles}.
A medida que se baja en la jerarquía, el dominio de datos que maneja cada interviniente es
cada vez más chico.

\begin{center}
\begin{tikzpicture}[font=\small,>=Stealth,
  role/.style={draw,rounded corners=2pt,minimum height=0.95cm,align=center,
               text=white,font=\bfseries},
  note/.style={align=left,font=\footnotesize,text width=4.4cm}]

  \node[role,fill=attackRed,minimum width=4.0cm] (trustee) {Data Trustee};
  \node[role,fill=orange!85!black,minimum width=2.2cm,below left=1.0cm and 0.1cm of trustee] (s1) {Data Steward};
  \node[role,fill=orange!70!yellow,minimum width=2.2cm,below right=1.0cm and 0.1cm of trustee] (s2) {Data Steward};
  \node[role,fill=orange!60!yellow,minimum width=1.7cm,below=1.0cm of s1] (c1) {Data Custodian};
  \node[role,fill=orange!75!black,minimum width=1.7cm,below=1.0cm of s2] (c2) {Data Custodian};
  \node[role,fill=itbaCyanDark,minimum width=5.2cm] (user) at ($(c1)!0.5!(c2) - (0,1.5cm)$) {Data User};

  \draw (trustee) -- (s1);
  \draw (trustee) -- (s2);
  \draw (s1) -- (c1);
  \draw (s2) -- (c2);
  \draw[<->] (c1) -- (user);
  \draw[<->] (c2) -- (user);
\end{tikzpicture}
\end{center}

\begin{description}[leftmargin=0pt,itemsep=4pt]
  \item[\textcolor{attackRed}{Data Trustee.}] Posición estratégica alta (nivel \emph{cabinet}),
        \textbf{responsable legalmente} por los datos. Si pasa algo con esos datos, es el
        responsable a nivel jurídico.
  \item[\textcolor{itbaCyanDark}{Data Steward.}] Responsable interno que \textbf{aprueba
        accesos} a los datos sensibles/restringidos \emph{dentro de su dominio}.
  \item[\textcolor{itbaCyanDark}{Data Custodian.}] Autorizado por un Steward para
        \textbf{otorgar (grant) acceso} a los usuarios sobre elementos del dominio del Steward.
  \item[\textcolor{itbaCyanDark}{Data User.}] El individuo autorizado que efectivamente
        \textbf{accede a los datos} y hace algo con ellos.
\end{description}

%=====================================================================
\section{Descubrimiento (Discovery)}
%=====================================================================

Identificar y «contar» los datos que se van a proteger. No todo es un activo y no todo dato
requiere asegurarse, así que primero hay que \textbf{clasificar} y luego decidir.

\subsection{Clasificación de los datos}

La diapositiva la presenta como un ciclo \textbf{PLAN $\to$ DO $\to$ CHECK $\to$ ACT}:

\begin{center}
\begin{tikzpicture}[font=\small,>=Stealth,
  phase/.style={draw,circle,fill=practGreen,text=white,font=\bfseries,minimum size=1.0cm},
  lbl/.style={align=center,font=\footnotesize}]

  \def\R{2.6cm}
  \node[phase] (p) at (90:\R)  {PLAN};
  \node[phase] (d) at (0:\R)   {DO};
  \node[phase] (c) at (-90:\R) {CHECK};
  \node[phase] (a) at (180:\R) {ACT};

  \node[lbl,above=2pt of p] {Identificar activos · roles · perfiles};
  \node[lbl,right=2pt of d,text width=3cm] {Educar · desplegar tecnología};
  \node[lbl,below=2pt of c] {Revisar logs y alertas};
  \node[lbl,left=2pt of a,text width=2.6cm] {Reclasificar · bloquear/permitir};

  \draw[->,line width=1.2pt,practGreen] (p) to[bend left=20] (d);
  \draw[->,line width=1.2pt,practGreen] (d) to[bend left=20] (c);
  \draw[->,line width=1.2pt,practGreen] (c) to[bend left=20] (a);
  \draw[->,line width=1.2pt,practGreen] (a) to[bend left=20] (p);
\end{tikzpicture}
\end{center}

\begin{itemize}[leftmargin=1.4em,itemsep=2pt]
  \item \textbf{Planificar:} identificar los \emph{activos}, identificar los \emph{roles}
        sobre esos datos (quién controla cada dominio — los \emph{data roles}) y definir
        \emph{perfiles de acceso} (algunos no requieren autorización; otros más protegidos).
  \item \textbf{Ejecutar:} \emph{educar} a todos los involucrados sobre la clasificación de
        los datos y quién tiene potestad sobre cada uno; \emph{desplegar la tecnología} que
        permita protegerlos (herramientas, políticas y perfiles de acceso).
  \item \textbf{Monitorear:} revisar \emph{logs y alertas} para verificar que no se vulnere
        el acceso a datos que no corresponden.
  \item \textbf{Reaccionar:} si se detectan desvíos, \emph{reclasificar} y \emph{bloquear o
        permitir} según corresponda (un dato crítico puede pasar a no crítico y viceversa).
\end{itemize}

\subsection{El mismo dato en distintos estados}

El mismo dato puede pasar por distintos \textbf{estados}, formatos, medios y tecnologías, en
distintas ubicaciones (hoy los data centers están distribuidos por el mundo). Las regulaciones
definen lineamientos para \textbf{cada estado por separado}. Habitualmente se hablaba de dos
estados (\emph{en reposo} y \emph{en tránsito}), pero hay un tercero importante (\emph{en uso}).

\begin{defbox}{Los tres estados del dato}
\begin{description}[leftmargin=0pt,itemsep=2pt,topsep=2pt]
  \item[\textbf{Data at-rest (en reposo).}] El dato está quieto, no se mueve ni se procesa.
        Ej.: una base de datos, un backup, una hoja de Excel, un objeto en un bucket S3.
  \item[\textbf{Data in-motion (en movimiento/tránsito).}] El dato se mueve de una ubicación o
        estado a otro, generalmente por la red (también un pendrive, o cuando una app le pide
        un dato a otra y el dato sale de la base hacia el programa).
  \item[\textbf{Data in-use (en uso).}] El dato está siendo procesado/utilizado por una
        aplicación, un servidor o una persona (transformándolo, calculando algo con él, etc.).
\end{description}
\end{defbox}

\noindent Sobre el mismo flujo de datos del problema inicial, cada componente está en alguno de
los tres estados:

\begin{center}
\begin{tikzpicture}[
  font=\footnotesize,>=Stealth, node distance=0.85cm and 0.85cm,
  store/.style={cylinder,shape border rotate=90,draw,fill=stateUse,minimum width=1.4cm,minimum height=0.85cm,align=center},
  app/.style={rounded corners=6pt,draw,fill=stateUse,minimum width=1.5cm,minimum height=0.95cm,align=center},
  doc/.style={draw,signal,signal to=south,fill=stateRest,minimum width=1.4cm,minimum height=0.85cm,align=center,inner sep=2pt},
  cl/.style={cloud,draw,aspect=2,minimum width=1.6cm,minimum height=1.0cm,align=center},
  arr/.style={->,line width=2.2pt,stateMotion}]

  \node[store] (rh) {RR.HH.};
  \node[app, right=of rh] (app) {RR.HH.\\APP};
  \node[cl, fill=stateUse, below right=0.85cm and 0.1cm of app] (drive) {Drive};
  \node[app, right=1.6cm of app] (xls) {Excel};
  \node[app, right=of xls] (mail) {Email};
  \node[cl, fill=white, right=of mail] (nube) {};

  \node[doc, below=0.5cm of rh] (b1) {BACKUP};
  \node[doc, below=0.4cm of drive] (b2) {BACKUP};

  \draw[arr] (rh) -- (app);
  \draw[arr] (rh) -- (b1);
  \draw[arr] (app) |- (drive);
  \draw[arr] (drive) -- (b2);
  \draw[arr] (app.east) -- ([yshift=4pt]xls.west);
  \draw[arr] (drive.east) -| (xls.south);
  \draw[arr] (xls) -- (mail);
  \draw[arr] (mail) -- (nube);

  % leyenda
  \node[draw,fill=stateRest,rounded corners=2pt,below=1.3cm of rh,anchor=west,xshift=-0.4cm] (lr)
       {\textbf{Data at-rest}};
  \node[draw,fill=stateMotion,rounded corners=2pt,right=0.5cm of lr] (lm) {\textbf{Data in-motion}};
  \node[draw,fill=stateUse,rounded corners=2pt,right=0.5cm of lm] (lu) {\textbf{Data in-use}};
\end{tikzpicture}
\end{center}

\subsection{Data Inventory}

\begin{defbox}{Data Inventory}
Identifica y \textbf{registra todos los datos críticos}: su ubicación, su uso, su
clasificación y los responsables de cada uno. Sirve de \textbf{guía de referencia} para la
toma de decisiones, el diseño de la seguridad y la aplicación de políticas.
\end{defbox}

No alcanza con «proteger todos los datos y listo»: los datos se necesitan para procesamiento,
análisis y para que distintas personas los usen en distintos formatos. Por eso conviene
inventariarlos: saber que existen, cómo deben clasificarse y de qué forma se pueden usar. El
ciclo de inventario es: \emph{Engage the business $\to$ Collect the information $\to$ Build the
inventory $\to$ Create classifications $\to$ Use the inventory}.

\subsubsection*{Desafíos del Data Inventory}

\begin{itemize}[leftmargin=1.4em,itemsep=2pt]
  \item \textbf{Costo de mantenerlo actualizado:} cambian los tipos de datos y, sobre todo,
        los \textbf{flujos de los datos} (la slide lo resalta en negrita). Cada vez que aparece,
        cambia o desaparece un dato, hay que actualizarlo.
  \item \textbf{Costo de inventariar:} descubrir nuevos tipos de datos es caro, especialmente
        en organizaciones muy grandes (hay que barrer cómo opera toda la empresa).
  \item \textbf{Falta de educación en su uso:} si la gente no avisa que generó datos nuevos,
        no se le saca provecho al inventario.
  \item \textbf{Son datos sensibles:} el propio inventario debe protegerse bajo las mismas
        reglas.
  \item \textbf{Complejo para la dinámica de las nuevas empresas:} queda rápidamente
        desactualizado.
\end{itemize}

\subsection{Data Labeling (etiquetado)}

Como alternativa al inventario centralizado surgió el \textbf{data labeling}.

\begin{defbox}{Data Labeling}
Busca inventariar los datos \textbf{dentro de los datos mismos} (como \emph{metadata}):
el dato lleva su propia clasificación. \textbf{Descentraliza} el inventario, pero permite que
las etiquetas \textbf{circulen con los datos}. Simplifica las reglas de protección y monitoreo.
\end{defbox}

\begin{itemize}[leftmargin=1.4em,itemsep=2pt]
  \item Es \textbf{más flexible y realizable} que el inventario tradicional: el sistema que
        genera el dato ya le asigna su etiqueta.
  \item \textbf{No es un proceso estandarizado:} cada sistema debe \emph{traducir} las
        etiquetas de un medio a otro. La organización puede normalizarlo internamente («todos
        los que manejen este tipo de datos sensibles aplican la política del labeling»).
  \item La mayoría de los sistemas modernos permiten crear etiquetas/metadata donde guardar
        esta información.
\end{itemize}

\subsubsection*{Ejemplo: AWS Tag Policy}

Una política de IAM de AWS otorga acceso a un objeto S3 sólo si el objeto tiene cierta
etiqueta. Es la materialización de la idea de \emph{control de acceso basado en etiquetas}:

\begin{lstlisting}[style=json]
{
  "Version": "2012-10-17",
  "Statement": [
    {
      "Principal": {                                  // <-- SUJETO
        "AWS": [ "arn:aws:iam::111122223333:role/JohnDoe" ]
      },
      "Effect": "Allow",
      "Action": ["s3:GetObject", "s3:GetObjectVersion"],
      "Resource": "arn:aws:s3:::DOC-EXAMPLE-BUCKET/*", // <-- OBJETOS
      "Condition": {
        "StringEquals": {
          "s3:ExistingObjectTag/environment": "production" // <-- ETIQUETA
        }
      }
    }
  ]
}
\end{lstlisting}

\begin{itemize}[leftmargin=1.4em,itemsep=1pt]
  \item \textbf{Sujeto:} el \texttt{Principal} (el rol \texttt{JohnDoe}).
  \item \textbf{Objetos:} el \texttt{Resource} (los objetos del bucket).
  \item \textbf{Etiqueta:} la \texttt{Condition} sobre la tag \texttt{environment = production}.
\end{itemize}

%=====================================================================
\section{Protección (Protect)}
%=====================================================================

Toda la preparación anterior (gobierno, descubrimiento, clasificación) sirve para
\textbf{efectivamente proteger} los datos. Acá se ven las técnicas concretas.

\subsection{Enforcement en tiempo de compilación (Compiler-time)}

\begin{defbox}{Enforcement (Compiler-time)}
Asegura que el código cumple ciertas políticas de seguridad \textbf{antes de ser ejecutado}.
Se hace en la fase de compilación, donde el compilador analiza el código fuente para detectar
posibles violaciones (\textbf{análisis estático de código}). El compilador puede asegurar que
los datos clasificados como confidenciales no se filtren a variables o zonas de memoria de
menor seguridad.
\end{defbox}

Ejemplos de violaciones que se buscan detectar: imprimir contraseñas o datos sensibles en un
log. El profesor aclara una errata de la slide: el ejemplo en Java es el
\texttt{GuardedString} (la slide dice \texttt{SecureString}).

\begin{notebox}{\texttt{GuardedString} vs.\ \texttt{String} en Java}
Un \texttt{GuardedString} guarda el string \textbf{encriptado}: si alguien inspecciona la
memoria en vivo, no lo va a poder leer (tendría que desencriptarlo). En cambio, una contraseña
en un \texttt{String} clásico de Java se ve en \textbf{texto plano} examinando la memoria del
proceso con una herramienta. El profe acota que el \texttt{GuardedString} es «bastante lento».

\smallskip
También ayudan las \textbf{herramientas de análisis de código estático}, que conectan con las
metodologías \emph{DevSecOps} vistas la clase anterior (meter la seguridad dentro del
desarrollo).
\end{notebox}

\subsection{Encriptación de datos en reposo (at-rest)}

Los datos en reposo no se mueven; una de las medidas es \textbf{encriptarlos en el medio en
que están}. Como esos datos después hay que recuperarlos, se usa encriptación (reversible).
Aparecen varias alternativas de \emph{dónde} se guardan las claves:

\begin{description}[leftmargin=0pt,itemsep=4pt]
  \item[\textbf{HSM} (\emph{Hardware Security Module}, FIPS 140-1 a 4).] Hardware específico de
        seguridad (visto también en Autenticación) dedicado a proteger secretos —p.\ ej.\ las
        \textbf{claves simétricas}, que son muy sensibles porque con una sola se puede
        desencriptar un montón de datos. Es \textbf{anti-tamper}: no se puede vulnerar
        fácilmente. Además, \textbf{ofrece las operaciones de (des)encriptación dentro del
        propio hardware}, así uno nunca accede a la clave: la operación la hace el HSM.
  \item[\textbf{Keystore.}] Forma de resguardar claves (con software asociado): se guardan
        todas las claves encriptadas con una \textbf{clave maestra}, y con esa clave se accede
        a las demás.
  \item[\textbf{Security Appliance.}] Otro dispositivo dedicado a este resguardo (de la slide).
\end{description}

\subsubsection*{Esquema KEK-DEK (encriptación por capas)}

\begin{defbox}{KEK / DEK}
\textbf{DEK} = \emph{Data Encryption Key} (la clave que efectivamente cifra los datos).\\
\textbf{KEK} = \emph{Key Encryption Key} (una clave que cifra a la DEK).
\end{defbox}

\begin{center}
\begin{tikzpicture}[font=\small,>=Stealth,node distance=0.9cm and 1.4cm,
  blk/.style={draw,rounded corners=2pt,minimum height=0.9cm,minimum width=1.5cm,align=center},
  op/.style={draw,fill=black!75,text=white,rounded corners=2pt,minimum size=0.85cm},
  key/.style={draw,fill=itbaBlue!40,rounded corners=2pt,minimum height=0.8cm,minimum width=1.2cm,align=center}]

  \node[blk] (data) {Data};
  \node[op,right=of data] (ed1) {E/D};
  \node[cylinder,shape border rotate=90,draw,right=of ed1,minimum width=1.3cm,minimum height=0.9cm] (db) {DB};
  \node[key,below=of ed1] (dek) {DEK};
  \node[op,below right=0.9cm and 0.9cm of dek] (ed2) {E/D};
  \node[key,below=of ed2] (kek) {KEK};

  \draw[->,line width=1.2pt] (data) -- (ed1);
  \draw[->,line width=1.2pt] (ed1) -- (db);
  \draw[->,line width=1pt] (dek) -- (ed1);
  \draw[->,line width=1pt] (kek) -- (ed2);
  \draw[->,line width=1pt] (ed2) -- (dek);
\end{tikzpicture}
\end{center}

\begin{itemize}[leftmargin=1.4em,itemsep=2pt]
  \item \textbf{Más seguro:} con la KEK se cifra la DEK, que es la que termina cifrando los
        datos (dos capas).
  \item \textbf{Más flexible:} se puede \textbf{rotar la KEK} sin tener que reencriptar todos
        los datos, porque la DEK (la que cifró los datos) no cambia.
  \item \textbf{Contra:} es \textbf{más difícil de administrar} (implica una segunda capa de
        encriptación).
  \item Ejemplos de cifrado de almacenamiento en reposo: \textbf{EncFS / Loop-AES},
        \textbf{EFS} (Windows), \textbf{FileVault} (macOS).
  \item La slide marca como problemas pendientes (\(\times\)): \textbf{Data Remanence}
        (remanencia de datos) y \textbf{Secure Deletion} (borrado seguro).
\end{itemize}

\subsubsection*{En reposo, en la nube}

La nube es \textbf{infraestructura de terceros}, en posesión del proveedor. La estrategia de
encriptación es un \emph{trade-off}: a más control sobre las claves, mayor complejidad y menor
costo del lado del proveedor; a menos control, más eficiencia operativa pero mayor costo.

\begin{center}
\begin{tikzpicture}[font=\footnotesize,>=Stealth,
  opt/.style={draw,rounded corners=2pt,minimum height=1.1cm,text width=2.0cm,align=center}]

  \node[opt,fill=attackRed!18] (smk) {Service\\Managed Keys};
  \node[opt,fill=attackRed!35,right=0.4cm of smk] (cmk) {Customer\\Managed Keys};
  \node[opt,fill=attackRed!55,right=0.4cm of cmk] (byok) {Bring Your\\Own Keys};
  \node[opt,fill=attackRed!80,text=white,right=0.4cm of byok] (hyok) {Hold Your\\Own Keys};

  \draw[<-,line width=3pt,itbaCyan]
    ([yshift=-12pt]smk.south west) -- node[below,black,font=\scriptsize]{Eficiencia operativa} ([yshift=-12pt]hyok.south east);
  \draw[->,line width=3pt,itbaCyanDark]
    ([yshift=-26pt]smk.south west) -- node[below,black,font=\scriptsize]{Seguridad / control propio} ([yshift=-26pt]hyok.south east);
\end{tikzpicture}
\end{center}

\begin{center}
\renewcommand{\arraystretch}{1.3}
\begin{tabular}{>{\bfseries}p{3.6cm} p{11.0cm}}
\toprule
\multicolumn{2}{l}{\itshape Cloud Workload con encriptación \emph{cloud-native} (claves del lado del proveedor):}\\
\midrule
Service Managed Keys & El cliente \textbf{no gestiona} las claves. El proveedor decide cuándo
rotarlas. No se puede acceder a las claves privadas ni usarlas para cifrar otros contenidos
distintos de los que el proveedor decide. (Máxima delegación, máximo costo.) \\
Customer Managed Keys & El cliente \textbf{gestiona} las claves y decide cuándo rotarlas; no
puede acceder a las claves privadas, pero puede encriptar lo que desee. \\
Bring Your Own Keys & El cliente \textbf{trae sus propias claves} (las genera y se las da al
proveedor), las posee y decide cuándo rotarlas. \\
\midrule
\multicolumn{2}{l}{\itshape Cloud Workload con encriptación propia:}\\
\midrule
Hold Your Own Keys & El cliente cifra \textbf{de su lado, antes de almacenar} la info en el
proveedor. Independiente del proveedor de nube (máximo control, mínima delegación). \\
\bottomrule
\end{tabular}
\end{center}

\begin{notebox}{¿De qué depende la elección? (pregunta de M.\ Gallardo en clase)}
Depende 100\% de tres cosas: \textbf{(1)} qué querés hacer / cuánto te importa la seguridad;
\textbf{(2)} cuánto querés delegar; y \textbf{(3)} tu posibilidad económica. Cuanto más a la
izquierda (\emph{Service Managed}), más responsabilidad tiene el proveedor $\Rightarrow$ mayor
costo, menos complejidad para vos. Cuanto más a la derecha (\emph{Hold Your Own}), menos costo
pero más complejidad propia (rotar y generar tus claves, protegerlas). Muchas empresas no son
expertas en seguridad y terminan delegando cada vez más.
\end{notebox}

\begin{warnbox}{Ojo con la ubicación de los datos en la nube}
Los proveedores tienen distintas \textbf{regiones}. Según cómo esté armada la distribución,
los datos pueden quedar repartidos en \textbf{distintos países con regulaciones distintas}
cada uno. Hay que tenerlo muy en cuenta.
\end{warnbox}

\subsection{Encriptación de datos en tránsito (in-transit)}

Los datos en tránsito se mueven de una ubicación a otra (internet, red privada). Dos enfoques:

\begin{enumerate}[leftmargin=1.6em,itemsep=2pt]
  \item \textbf{Encriptar antes de mover el dato.} Ej.: cifrar un archivo con \textbf{GPG}
        antes de mandarlo por mail.
  \item \textbf{Usar protocolos de transporte que cifran.} Ej.: enviar el mail con
        \textbf{SMTP sobre TLS}; así no hace falta cifrar el archivo a mano y igual viaja
        seguro.
\end{enumerate}

\noindent Protocolos de cifrado in-transit: \textbf{TLS}, \textbf{IPSEC}, \textbf{SSH}. Todo el
tráfico \textbf{HTTPS} va asegurado por TLS (no se puede leer). El esquema general es el ya
visto: se acuerda primero con \textbf{criptografía asimétrica} la \textbf{clave simétrica}, y
luego se cifra todo el tráfico de forma simétrica.

\subsection{Enforcement en tiempo de ejecución (Execution-time)}

\begin{defbox}{Enforcement (Execution-time)}
Es la técnica más utilizada para \textbf{monitorear y controlar cómo se transmite la
información entre los componentes} de un sistema. Requiere entender el \textbf{movimiento de
datos} dentro del sistema de información: entre procesos, entre distintos niveles de seguridad
o entre distintos estados. Se implementa:
\begin{itemize}[leftmargin=1.4em,topsep=2pt,itemsep=1pt]
  \item a \textbf{nivel aplicativo} (buenas decisiones y principios de seguridad: no imprimir
        secretos en un log, tomar medidas para no exponer datos sensibles), o
  \item a \textbf{nivel de infraestructura}, con herramientas \textbf{DLP} (\emph{Data Loss
        Prevention}) que detectan filtraciones/exposición de datos.
\end{itemize}
\end{defbox}

\subsection{Encriptación en memoria}

Los datos en memoria son un caso particular: se los puede ver como \emph{at-rest} en memoria
esperando que la CPU los use, o \emph{in-transit} entre el bus de I/O y la CPU. El riesgo:
en un ambiente de \textbf{múltiples máquinas virtuales accediendo a la misma memoria}, hay
riesgo potencial de acceder a información confidencial de otra VM.

\begin{defbox}{TME y TME-MK (Intel y otros fabricantes)}
\textbf{TME} = \emph{Total Memory Encryption}: mecanismo implementado en \textbf{hardware}, con
una \textbf{sola clave} de encriptación compartida por todos los procesos. La desencriptación
se realiza \textbf{dentro de la CPU, antes de la caché}.\\
\textbf{TME-MK} = \emph{Total Memory Encryption — Multi-Key}: \textbf{múltiples claves} de
encriptación, con un control de acceso soportado por los registros de \textbf{VT-x}. El cifrado
es \emph{AES-XTS} y cada página de memoria física lleva asociado un \emph{KeyID}.
\end{defbox}

\begin{center}
\begin{tikzpicture}[font=\footnotesize,>=Stealth,
  box/.style={draw,fill=black!10,minimum width=2.0cm,minimum height=0.6cm,align=center},
  ctrl/.style={draw,fill=black!55,text=white,minimum width=1.4cm,minimum height=0.55cm,align=center}]

  \node[box] (core) {Cores};
  \node[box,below=0.25cm of core,minimum width=3.0cm] (cache) {Cache};
  \node[ctrl,below left=0.3cm and -1.0cm of cache] (aes1) {AES-XTS};
  \node[ctrl,right=0.3cm of aes1] (aes2) {AES-XTS};
  \node[box,below=0.2cm of aes1,minimum width=1.4cm] (dram) {DRAM};
  \node[box,below=0.2cm of aes2,minimum width=1.4cm] (nvram) {NVRAM};
  \node[align=center,below=0.25cm of dram,text width=3.4cm,font=\scriptsize] (lbl)
       {Memoria cifrada con AES-XTS\\(la CPU descifra antes de la caché)};

  \draw[->] (dram) -- (aes1);
  \draw[->] (nvram) -- (aes2);
  \draw[->] (aes1) -- ([xshift=-0.6cm]cache.south);
  \draw[->] (aes2) -- ([xshift=0.6cm]cache.south);
  \draw[<->] (cache) -- (core);

  \node[draw,fill=stateUse,rounded corners=2pt,right=1.6cm of core,align=center] (vm1) {VM1\\KeyID 0,\textbf{1},3};
  \node[draw,fill=black!8,rounded corners=2pt,right=0.4cm of vm1,align=center] (vm2) {VM2\\KeyID 0,\textbf{2},3};
  \node[draw,below=0.4cm of vm1,xshift=0.9cm,align=center,text width=2.6cm,font=\scriptsize] (vmm)
       {VMM: usa KeyID 0 para sus páginas y gestiona los KeyID de toda la memoria};
\end{tikzpicture}
\end{center}

\subsection{Homomorphic Encryption}

\begin{defbox}{Encriptación homomórfica}
Es un esquema de encriptación que admite \textbf{operar sobre el dato cifrado}: aplicar una
función al texto cifrado equivale a aplicarla al texto plano. La slide lo expresa como:
\[
  f(x) = y \;\Longrightarrow\; f(\mathrm{Enc}(x)) = \mathrm{Enc}(f(x))
\]
\end{defbox}

\begin{warnbox}{Nota de fidelidad sobre la fórmula de la slide}
La diapositiva escribe literalmente \(f(x)=y \Rightarrow f(\mathrm{Enc}(x)) =
f(\mathrm{Dec}(y))\). La \emph{idea} que transmite —y que el profe describe como «un esquema
que admite operar»— es la propiedad homomórfica: poder calcular sobre datos cifrados sin
desencriptarlos. Se reproduce la fórmula tal cual aparece para no alterar la fuente, aclarando
que la formulación canónica del homomorfismo es \(f(\mathrm{Enc}(x)) = \mathrm{Enc}(f(x))\).
\textbf{Referencia citada en la slide:} \emph{Protecting privacy through homomorphic
encryption}, Lauter, 2022.
\end{warnbox}

\subsection{Encriptación en SoC / IoT / sistemas embebidos}

Uno de los escenarios \textbf{más desfavorables}: chips pequeños (Arduino, Raspberry,
ESP8266/ESP32) metidos en cualquier lado (un semáforo, un campo de cultivos, dentro de otras
cosas). Características y vectores de ataque:

\begin{itemize}[leftmargin=1.4em,itemsep=2pt]
  \item \textbf{Tampering-with:} están \textbf{físicamente expuestos}; alguien puede ir y
        romperlos o manipularlos.
  \item \textbf{Open Arch (arquitectura abierta):} usan protocolos abiertos, ya muy conocidos.
  \item \textbf{Debug complejo (debug/seguridad):} al ser de bajo nivel, suelen permitir
        \emph{debugging} fácil, muchas veces \textbf{sin interfaz de autenticación}.
  \item \textbf{Wireless:} aunque el dispositivo esté físicamente inaccesible, uno se puede
        conectar inalámbricamente (hasta desde un celular).
  \item \textbf{Diversidad / Escala:} se despliegan en grandes cantidades, distribuidos por
        todos lados y dentro de muchas cosas.
  \item \textbf{Riesgo físico.}
\end{itemize}

La slide ubica estos riesgos en tres capas de la taxonomía de IoT: \textbf{Perception Layer},
\textbf{Application Layer} y \textbf{Network Layer} (chip de ejemplo: \emph{Microchip
ATECC608A}).

\begin{notebox}{Ejemplo del profe: sensores de riego}
Si un atacante manipula (\emph{tampering}) un sensor de riego controlado por uno de estos
chips —por ejemplo, fijando una condición de temperatura que sabe que se va a superar todo el
año— logra que el sistema riegue todo el tiempo y genere \textbf{gastos económicos}. Por eso
es importante meter algún tipo de \textbf{HSM / seguridad embebida} en estos aparatos, algo a
lo que todavía no se le da la importancia que tiene.
\end{notebox}

%=====================================================================
\section{Cumplimiento (Comply)}
%=====================================================================

\begin{redslide}
«No basta con ser bueno, también hay que parecerlo.»
\end{redslide}
\noindent Con esa anécdota («mi abuelo era un genio», sobre la viñeta de Daniel Paz), el profe
introduce el cumplimiento: \textbf{no alcanza con proteger; hay que poder demostrarlo}.

Todas las regulaciones internacionales y comerciales definen reglas e imponen un proceso de
\textbf{auditoría y certificación}. Cuando llega una auditoría, hay que \textbf{demostrar con
evidencia} que todo lo que dijiste que hacías, lo hacés efectivamente. Por eso las empresas
suelen obligar a que \textbf{todo se loguee} y quede registro: con esos registros es mucho más
fácil probar que se es \emph{compliant}.

\subsection{Retención de datos}

\begin{defbox}{Retención de la información}
Es el \textbf{guardado de información que ya no es necesaria}, usualmente por temas
regulatorios. Tiene fines \textbf{legales} o de \textbf{análisis forense}.
\end{defbox}

\begin{center}
\renewcommand{\arraystretch}{1.3}
\begin{tabular}{>{\bfseries}p{2.6cm} p{11.6cm}}
\toprule
HIPAA & Si un log, nota o registro se relaciona con una política/procedimiento HIPAA, debe
retenerse \textbf{6 años} desde la fecha en que el contenido fue usado por última vez / fue
efectivo por última vez. (Ej.: una obra social o medicina prepaga conserva la historia clínica
6 años aunque la persona deje de ser cliente.) \\
PCI-DSS & El \emph{Requirement 3.1} exige retener y seguir procedimientos de retención y
disposición de datos: los registros que ya no se necesitan deben borrarse pronta y
adecuadamente. \\
SOX & Retención de \textbf{7 años}: toda la documentación de auditoría/revisión de los estados
financieros, por si hay que investigar a futuro. \\
\bottomrule
\end{tabular}
\end{center}

\begin{notebox}{Rotación de claves por cumplimiento}
Algunas leyes que obligan a encriptar ciertos datos personales también obligan a
\textbf{rotar las claves de encriptación} (p.\ ej.\ cada 10 años) y \textbf{reencriptar} todo.
Si a los 10 años llega una auditoría y no rotaste, probablemente no estés cumpliendo y tengas
una consecuencia (probablemente económica).
\end{notebox}

%=====================================================================
\section{Detección (Detect): Data Loss Prevention (DLP)}
%=====================================================================

Cumplir no es suficiente: hay que \textbf{verificar constantemente} que se sigue cumpliendo.
La herramienta principal de esta etapa son los \textbf{DLP}.

\begin{defbox}{DLP — Data Loss Prevention}
Herramientas que ayudan en la \textbf{clasificación, inventariado y detección} de los datos en
todos sus estados. A pesar del nombre («Prevention»), \textbf{la mayoría son de detección, no
de prevención}: están orientadas a un control \textbf{detectivo}, no preventivo ni correctivo.
Se activan cuando se produce una falla en la aplicación de una política.
\end{defbox}

\begin{itemize}[leftmargin=1.4em,itemsep=2pt]
  \item Permiten ver los datos en todos sus estados y detectar si están saliendo de estados
        que no corresponden, si están siendo filtrados, divulgados o \textbf{exfiltrados}.
  \item Se le da toda la información a la herramienta y ésta \textbf{alerta} cuando algo no se
        está cumpliendo.
  \item \textbf{No son una panacea} y son \textbf{súper complejas de implementar}: requieren un
        esfuerzo enorme de toda la compañía.
\end{itemize}

\subsection{DLP hoy en día}

Como eran tan difíciles de implementar, hoy las \textbf{funcionalidades de DLP ya vienen
incorporadas} en muchos de los productos de uso diario (muchas empresas que «hacían DLP»
fueron compradas e integradas a otras soluciones). Los DLP cubren múltiples canales de salida:
USB/CD/DVD, Print/Fax, Email, Webmail, mensajería instantánea, FTP, servidores, bases de datos,
etc.

\begin{notebox}{Ejemplo: ransomware}
Una funcionalidad de DLP puede detectar \textbf{encriptaciones no autorizadas} de datos, o que
un dato cambió o tiene un comportamiento extraño. Primero alerta; y según la sofisticación de
la herramienta, podría tomar medidas adicionales (p.\ ej.\ \textbf{inutilizar el dispositivo}
de un empleado cuya máquina fue encriptada por un ataque).
\end{notebox}

%=====================================================================
\section{Respuesta (Respond): respuesta a incidentes}
%=====================================================================

\begin{defbox}{Respuesta a incidentes}
Cómo opera la organización cuando detecta una filtración, un ataque o una vulneración de la
protección de los datos. Por lo general, hay \textbf{equipos dedicados} en las empresas a la
respuesta a incidentes, que ya tienen el protocolo incorporado y van indicando los pasos al
resto de las personas.
\end{defbox}

\subsection{Qué hacer en las primeras 24 horas}

Durante el incidente se trata de \textbf{guardar, documentar e identificar} todo lo que pasó,
y de \textbf{frenar la filtración}. La checklist de la slide:

\begin{center}
\begin{tikzpicture}[font=\footnotesize,node distance=0.35cm,
  it/.style={draw,fill=black!72,text=white,rounded corners=2pt,minimum width=2.55cm,
             minimum height=1.0cm,align=center,inner sep=3pt}]
  \node[it] (a) {Record the\\date and time};
  \node[it,right=of a] (b) {Alert\\everyone};
  \node[it,right=of b] (c) {Secure the\\premises};
  \node[it,right=of c] (d) {Stop data\\exfiltration};
  \node[it,right=of d] (e) {Document\\everything};

  \node[draw,fill=attackRed,text=white,below=0.35cm of c,minimum width=13.3cm,
        font=\bfseries,align=center] (mid) {WHAT TO DO IN THE FIRST 24 HOURS};

  \node[it,below=0.35cm of mid.west,anchor=north west] (f) {Interview\\anyone involved};
  \node[it,right=of f] (g) {Review comms\\protocols};
  \node[it,right=of g] (h) {Assess priorities\\and risks};
  \node[it,right=of h] (i) {Bring in a\\forensic team};
  \node[it,right=of i] (j) {Notify law\\enforcement};
\end{tikzpicture}
\end{center}

\noindent Se toman además medidas de \textbf{sobre-segurización} para evitar otras
exfiltraciones que no se estén viendo (la analogía del profe: si te entran a robar por una
puerta insegura, ponés una puerta súper segura y además asegurás ventanas, techo, etc.).

\begin{warnbox}{No ocultar el incidente — pero comunicar lo mínimo y necesario}
No hay que tratar de \textbf{ocultar} los incidentes: hay que trabajar con un protocolo. Pero
tampoco hay que sobre-exponer información: \emph{no} se sale a publicar en Twitter «ayer
hackearon la base de datos de mis clientes». Si se está sujeto a una normativa de datos
personales, hay que \textbf{informar al organismo} correspondiente, comunicando \textbf{sólo
lo requerido}. Después ese organismo pedirá la información que necesite, y ahí sí se le
disponibiliza (cómo fue el ataque, qué datos quedaron expuestos), tarea del \textbf{equipo de
forense} dentro de la respuesta a incidentes.
\end{warnbox}

\subsection{Más allá de las 24 horas}

\begin{center}
\begin{tikzpicture}[font=\footnotesize,>=Stealth,
  seg/.style={draw,rounded corners=2pt,minimum width=2.9cm,minimum height=0.9cm,align=center,text=white}]
  \node[align=center,font=\bfseries] (c) {BEYOND\\24 HOURS};
  \node[seg,fill=attackRed!75,above right=0.2cm and 0.6cm of c] (w) {Work with forensics};
  \node[seg,fill=itbaCyanDark,below right=0.2cm and 0.6cm of c] (l) {Identify legal obligations};
  \node[seg,fill=black!60,below=0.7cm of c] (r) {Report to upper management};
  \node[seg,fill=practGreen,below left=0.2cm and 0.6cm of c] (b) {Identify possible battles};
  \node[seg,fill=orange!85!black,above left=0.2cm and 0.6cm of c] (f) {Fix the issue};
\end{tikzpicture}
\end{center}

\begin{itemize}[leftmargin=1.4em,itemsep=2pt]
  \item \textbf{Trabajar con el equipo de forenses} e \textbf{identificar las implicancias
        legales} (qué hay que hacer legalmente). Ej.: ante una filtración de tarjetas de
        crédito, probablemente haya que informar a los \textbf{bancos} emisores qué pasó y qué
        clientes fueron afectados, y también a los propios clientes (el típico mail «tus datos
        fueron vulnerados»; en passwords, pedir el segundo factor o cambiar la contraseña).
  \item \textbf{Reportar a la dirección} (\emph{upper management}); si es una empresa pública,
        dar explicaciones al board / a los inversores.
  \item \textbf{Identificar posibles problemas} (\emph{battles}) que el incidente pueda traer a
        futuro.
  \item \textbf{Arreglar el problema} (\emph{fix the issue}): según la causa —si fue un
        empleado que divulgó algo, bajar una política exhaustiva; si fue una vulnerabilidad de
        un sistema, corregirla.
\end{itemize}

%=====================================================================
\section{Cierre y mapa mental}
%=====================================================================

Con esta clase \textbf{termina la materia}. El hilo conductor fue recorrer el ciclo de
protección de la información de punta a punta, aplicando sobre los datos los principios de la
\textbf{tríada CIA} vistos durante el curso.

\begin{center}
\begin{tikzpicture}[font=\footnotesize,>=Stealth,
  et/.style={draw,rounded corners=2pt,fill=itbaBlue!30,minimum width=2.6cm,
             minimum height=0.8cm,align=center}]
  \node[et] (g) {Govern\\\scriptsize Data Governance, roles, ciclo de vida};
  \node[et,below=0.3cm of g] (di) {Discover\\\scriptsize Clasificación, estados, inventory, labeling};
  \node[et,below=0.3cm of di] (pr) {Protect\\\scriptsize Encriptación at-rest / in-transit / in-use, DLP};
  \node[et,below=0.3cm of pr] (co) {Comply\\\scriptsize Auditoría, evidencia, retención, rotación};
  \node[et,below=0.3cm of co] (de) {Detect\\\scriptsize DLP (detectivo)};
  \node[et,below=0.3cm of de] (re) {Respond\\\scriptsize Forense, legal, primeras 24h y más allá};

  \foreach \a/\b in {g/di,di/pr,pr/co,co/de,de/re}
    \draw[->,line width=1pt] (\a) -- (\b);
\end{tikzpicture}
\end{center}

\begin{defbox}{Resumen de ideas para llevar}
\begin{itemize}[leftmargin=1.4em,itemsep=2pt,topsep=2pt]
  \item Los controles de acceso \textbf{no alcanzan}: el dato escapa del sistema; hay que
        protegerlo en todos sus estados y a lo largo de su ciclo de vida.
  \item \textbf{Acotar} los datos a proteger según criticidad (lección de PCI-DSS); usar
        \textbf{tokenización} y \textbf{anonimización} (PCI / HIPAA) cuando aplique.
  \item Tres estados del dato: \textbf{at-rest}, \textbf{in-motion}, \textbf{in-use}; cada uno
        con sus técnicas (HSM/keystore/KEK-DEK, TLS/IPSEC/SSH/GPG, TME-MK/homomórfica).
  \item En la nube, elegir el modelo de claves (\emph{Service / Customer / Bring / Hold Your
        Own Keys}) es un trade-off entre control, costo y complejidad; cuidado con la
        \textbf{ubicación legal} de los datos.
  \item \textbf{Comply:} «hay que parecerlo» — todo se loguea para tener \textbf{evidencia};
        respetar plazos de \textbf{retención}.
  \item \textbf{DLP} es detectivo, no preventivo, y complejo de implementar.
  \item Ante un incidente: documentar, contener, \textbf{no ocultar} pero comunicar lo
        mínimo y necesario, y trabajar con forense y legales.
\end{itemize}
\end{defbox}

\vfill
\begin{center}\small\itshape\color{black!60}
Apunte integral de la Clase 11 — Protección de datos · Criptografía y Seguridad, ITBA.\\
Síntesis fiel de la teoría (transcripción de la clase del Prof.\ Leandro Suárez + diapositivas oficiales).
\end{center}

\end{document}
